\chapter{Conclusions}
I recovered the plaintext of variant 19 through frequency analysis and
successive letter substitutions. The first seven recorded replacements
established THE, THESE, SYSTEM and THAT. Continuing with short words,
double letters and the emerging technical vocabulary allowed the remaining
observed letters to be identified. The full substitution was consistent
throughout the passage and reproduced the ciphertext when reversed.

The work demonstrates that letter frequencies are useful starting points,
but they do not determine the key on their own. For example, the third-most
frequent cipher symbol P represents S, while the rare symbol E represents
J. Digraphs, trigraphs, word patterns and sentence meaning supplied the
additional evidence needed to confirm these mappings.

I also learned to distinguish an observed substitution from one inferred
only by elimination: M does not appear in the ciphertext, so M$\to$Z
completes the assumed alphabet permutation without being directly tested
by the message. Preserving the exact input and decrypted versions makes
it possible to separate cryptographic correctness from corrections to
copied spacing and formatting.

Although monoalphabetic substitution has a large theoretical key space,
it preserves enough language structure to be vulnerable to this attack.
The main skills developed in this laboratory work were comparing
frequency distributions, testing hypotheses across several words,
maintaining a consistent substitution table and verifying a recovered
plaintext. No laboratory program was required.
